\documentclass{article}
\usepackage{amsmath}
\usepackage{graphicx}
\usepackage[letterpaper]{geometry}
\usepackage[utf8]{inputenc}
\usepackage[spanish]{babel}
\usepackage{gfsartemisia}
\usepackage[T1]{fontenc}

\title{Cálculo y simulación del campo magnético de una espira}
\author{}
\date{\today}

\begin{document}

\maketitle

Se parte de la definición de la ley de Biot-Savart, dada por:

\begin{equation}
	d \mathbf{B} = \frac{\mu_0 I}{4 \pi} \int \frac{ d \mathbf{s} \times \mathbf{r}' }{ | \mathbf{r}' |^3 }
\end{equation}

\section{Solución analítica}

Usamos la ecuación 1.

\section{Simulación}
% texto texto

\end{document}
